\section{On-Chain Credit Registry}

\label{sec:registry}
% ══════════════════════════════════════════════════════════════════════════════

\subsection{Credit Lifecycle}

Each carbon credit follows a lifecycle managed by smart contracts:

\begin{enumerate}
  \item \textbf{Project registration}: Project boundaries, methodology,
    and baseline parameters are recorded on-chain.
  \item \textbf{Monitoring}: CarbonNet biomass estimates and change
    detection results are committed as merkle roots at each monitoring
    interval.
  \item \textbf{Issuance}: Credits are minted as ERC-721 tokens when
    verified carbon stock gains exceed the dynamic baseline by a
    statistically significant margin.
  \item \textbf{Buffer pool}: 20\% of issued credits are held in a
    permanence buffer pool, released over the project crediting period
    if no reversals occur.
  \item \textbf{Retirement}: Credits are permanently burned when used
    for offsetting claims, with the retirement event linked to the
    retiree's account.
  \item \textbf{Reversal adjustment}: If reversals are detected, buffer
    pool credits are automatically cancelled proportional to the
    biomass loss.
\end{enumerate}

\subsection{Credit Metadata}

Each credit token contains or references:
\begin{itemize}
  \item Project ID and spatial boundary (GeoJSON hash).
  \item Vintage year and monitoring period.
  \item Biomass estimate with uncertainty bounds.
  \item Dynamic baseline value.
  \item Net sequestration (credit quantity in \tCO{}).
  \item Links to raw satellite imagery (IPFS CIDs).
  \item CarbonNet model version and parameters.
  \item Reversal risk score.
\end{itemize}

\subsection{Transparency and Auditability}

Any party can independently verify a credit by:
\begin{enumerate}
  \item Downloading the referenced satellite imagery from IPFS.
  \item Running the open-source CarbonNet model.
  \item Comparing the result against the on-chain biomass attestation.
  \item Verifying the dynamic baseline against published control areas.
\end{enumerate}

This creates a fully reproducible verification chain---a fundamental
departure from the opaque auditor reports of traditional carbon
standards.


% ══════════════════════════════════════════════════════════════════════════════
